% Subsection: Closure Under Complements  (notes stub; content authored later)
\subsection{Closure Under Complements}

\begin{remark*}[Exposition]
Complement closure is different from union or intersection closure because it
depends explicitly on the ambient set.  The complement of \(A\) is not
absolute; it is the part of the fixed universe \(X\) outside \(A\).
\end{remark*}

\begin{definition}[Closed Under Complements]\label{def:closed-under-complements}
Let \(X\) be a set and let \(\mathcal{F}\subseteq\mathcal{P}(X)\).  The family
\(\mathcal{F}\) is \emph{closed under complements} if
\[
  A\in\mathcal{F}
  \quad\Longrightarrow\quad
  X\setminus A\in\mathcal{F}.
\]
\end{definition}

\begin{remark*}[Standard quantified statement]
\[
  \forall A\in\mathcal{F},\qquad A^c=X\setminus A\in\mathcal{F}.
\]
\end{remark*}

\begin{remark*}[Predicate reading]
\[
\begin{aligned}
&C_{{}^c}=\mathsf{ClosureContext}(X,\mathcal{F},\mathcal{F},{}^c_X),\\
&\operatorname{ClosedUnder}({}^c_X,\mathcal{F},C_{{}^c}).
\end{aligned}
\]
\end{remark*}

\begin{remark*}[Interpretation]
Complement closure says that whenever a subset is admitted, its ambient
opposite is admitted too.  This makes the family stable under switching a
membership condition to its negation.
\end{remark*}

\begin{remark*}[Examples]
The full power set \(\mathcal{P}(X)\) is closed under complements.  The family
\(\{\varnothing,X\}\) is also closed under complements.  The family of finite
subsets of an infinite set \(X\) is not closed under complements, because the
complement of a finite subset of \(X\) may be infinite.
\end{remark*}

\begin{dependencies}
\begin{itemize}
  \item \hyperref[def:closed-under-operation]{Closed Under an Operation}
  \item \hyperref[def:complement]{Complement}
  \item \hyperref[def:power-set]{Power Set}
\end{itemize}
\end{dependencies}
